\chapter{Introduction}

\section{What an AMR Is}

An Autonomous Mobile Robot (AMR) navigates without fixed infrastructure (rails, magnetic strips). It uses LiDAR, cameras, and SLAM algorithms for real-time localization and collision avoidance. This distinguishes AMRs from Automated Guided Vehicles (AGVs), which follow predefined paths. Over the past decade, AMRs have moved from labs to commercial deployment in manufacturing, warehousing, healthcare, hospitality, and security \cite{standardbots_amr,valuates_agv_amr}.

\section{Why the Market Is Growing}

Three factors drive adoption:
\begin{itemize}
    \item \textbf{Labor shortages and rising wages} in mature economies make automation attractive.
    \item \textbf{Falling sensor and compute costs} (LiDAR, GPUs) have reduced hardware prices by 60--70\% over five years.
    \item \textbf{Maturing open-source software} (ROS/ROS2, Nav2) has lowered navigation development barriers.
\end{itemize}

A second trend is \textbf{modularity}: manufacturers now offer a common mobile base that accepts interchangeable payloads (delivery boxes, arms, sensor packs) instead of purpose-built robots \cite{fdatabot_amr2026,robotnik_kairos}. This report focuses on this modular segment.

\section{Scope and Method}

This study does \emph{not} assume international AMR success transfers directly to Tunisia. We test that assumption against local economic parameters: labor costs, company size, import dependency, and procurement practices. Global growth figures are used as context, not as forecasts.

The report covers market sectors (Chapters 2--3), competitors (Chapter 4), target clients, pricing, pain points, use cases (Chapters 5--8), and concludes with opportunities and recommendations (Chapters 9--10).